\documentclass[11pt]{article}
\usepackage[utf8]{inputenc}
\usepackage[margin=1in]{geometry}
\usepackage{amsmath}
\usepackage{graphicx}
\usepackage{booktabs}
\usepackage{hyperref}
\usepackage[sort&compress,numbers]{natbib}
\usepackage{float}

\title{A hippocampus--accumbens code guides goal-directed appetitive behavior}
\author{%
  Author~A$^{1,\ast}$, Author~B$^{1}$, Author~C$^{2}$, Author~D$^{1}$, Author~E$^{1,2,\ast}$\\[6pt]
  \small $^{1}$Institute of Neuroscience, Technical University of Munich, Munich, Germany\\
  \small $^{2}$Munich Cluster for Systems Neurology (SyNergy), Munich, Germany\\
  \small $^{\ast}$Correspondence: \texttt{corresponding@tum.de}
}
\date{}

\begin{document}
\maketitle

% ============================================================
% ABSTRACT
% ============================================================
\begin{abstract}
The hippocampus is critical for spatial memory and navigation, yet how specific hippocampal output pathways transmit information to guide goal-directed behavior remains poorly understood. Here we use dual-color two-photon calcium imaging in head-fixed Thy1-GCaMP6s mice to simultaneously record 5,372 dorsal hippocampal (dHPC) neurons, including 444 retrogradely labeled nucleus accumbens (NAc)-projecting neurons, across 19 sessions in 6 mice during a spatial appetitive task on a 360\,cm linear belt. We find that dHPC-to-NAc neurons are enriched for place cells (38\% vs.\ 32\%; $\chi^{2}(1, 5372)=6.364$, $p=0.012$) and exhibit higher spatial information rate (Welch's $t(186.55)=4.770$, $p<0.001$), sparsity ($t(218.79)=3.657$, $p<0.001$), reliability ($t(203.46)=2.479$, $p=0.014$), and in-field activity ($t(190.32)=2.884$, $p=0.0044$). Speed-inhibited neurons are overrepresented in the projection population (21\% vs.\ 15\%; $\chi^{2}(1, 5372)=13.66$, $p=0.00022$), as are appetitive lick-excited neurons (7.4\% vs.\ 3.8\%; $\chi^{2}(1, 5372)=13.018$, $p=0.00031$). Optogenetic activation of dHPC terminals in NAc evokes orofacial movement ($t(3)=7.485$, $p=0.00494$) and deceleration ($t(3)=-3.551$, $p=0.0381$). A generalized linear model confirms enhanced conjunctive coding for position, velocity, and licking in dHPC-to-NAc neurons (44\% vs.\ 19\% conjunctive; $p<0.001$), and conjunctive neurons decode reward zone presence more accurately than non-conjunctive neurons (Wilcoxon's $W(18)=14.0$, $p<0.001$). These results reveal that the hippocampus routes a multimodal, behaviorally relevant spatial code to the nucleus accumbens to guide appetitive goal-directed navigation.
\end{abstract}

\section{Introduction}

The hippocampus is a cornerstone of the mammalian memory and navigation systems, encoding spatial position through place cells that fire at specific locations within an environment~\citep{okeefe1971, moser2008}. Beyond purely spatial signals, hippocampal neurons integrate contextual, temporal, and reward-related information to form rich cognitive representations that support goal-directed behavior~\citep{eichenbaum2014, markus1995}. However, the hippocampus does not act in isolation; it routes its computations through multiple output pathways to downstream structures that translate mnemonic signals into adaptive actions~\citep{pennartz2011, wikenheiser2015}.

A key recipient of hippocampal output is the nucleus accumbens (NAc), a ventral striatal structure long hypothesized to serve as the interface between limbic and motor circuits~\citep{mogenson1980, floresco2015}. The hippocampal--accumbal pathway has been implicated in spatial context conditioning, conditioned place preference, and reward-based navigation~\citep{ito2008, leshan2014, trouche2019}. Disruption of this pathway impairs the ability to use spatial memories for goal-directed approach behavior, suggesting that the information transmitted from hippocampus to NAc is essential for translating spatial knowledge into motivated action~\citep{britt2012}.

Despite the established importance of this projection, the specific content of the hippocampal information stream reaching the NAc remains largely unknown. Hippocampal principal neurons are structurally and functionally diverse in terms of morphology, electrophysiology, and projection targets~\citep{danielson2016, grosmark2016}. Electrophysiological studies of projection-specific hippocampal coding have been limited by relatively low sample sizes or indirect connectivity measurements~\citep{ciocchi2015}. While recent advances in retrograde labeling and two-photon imaging have enabled projection-specific recordings~\citep{tervo2016, dombeck2010}, comprehensive characterization of the spatial, velocity, and appetitive coding properties of large populations of identified hippocampus-to-NAc projection neurons during behavior has remained elusive.

Here we address this gap by combining dual-color two-photon calcium imaging with retrograde viral labeling to simultaneously record from thousands of dorsal hippocampal neurons while identifying NAc-projecting cells in mice performing a spatial appetitive navigation task. We recorded a total of 5,372 GCaMP-expressing neurons including 444 putative dHPC-to-NAc neurons across 6 Thy1-GCaMP6s mice and 19 imaging sessions. Our approach enables high-throughput, projection-specific characterization of spatial coding, velocity modulation, appetitive licking signals, and conjunctive multimodal representations. We further use optogenetic activation of dHPC terminals in NAc to establish a causal link between this pathway and appetitive motor output. Our results reveal that the hippocampus-to-accumbens projection transmits an enriched, multimodal code that combines spatial, locomotor, and reward-related information to guide goal-directed appetitive behavior.

% ============================================================
% RESULTS
% ============================================================
\section{Results}

\subsection{Behavioral training on a spatial appetitive task}

Mice ($n=18$, two cohorts of 9) were trained over 5 days to navigate a 360\,cm textile belt with six textured zones and a 30\,cm hidden reward zone (Fig.~1a). Behavioral performance improved significantly across training days. Repeated-measures ANOVA revealed significant increases in rewarded laps per session ($F(4)=6.344$, Greenhouse--Geisser corrected, $p<0.05$), anticipatory licking ($F(4)=3.803$, $p<0.05$), and reward-zone licking ($F(4)=4.276$, Greenhouse--Geisser corrected, $p<0.05$; Table~\ref{tab:training}).

\begin{table}[H]
\centering
\caption{Training progression across 5 days assessed by repeated-measures ANOVA ($n=18$ mice).
$\uparrow$ indicates metrics expected to increase with learning.}
\label{tab:training}
\begin{tabular}{lccccc}
\toprule
Measure ($\uparrow$) & $F$-statistic & df & Correction & $p$ \\
\midrule
Rewarded laps per session & \textbf{6.344} & 4 & GG & $<0.05$ \\
Anticipatory licking & \textbf{3.803} & 4 & None & $<0.05$ \\
Reward-zone licking & \textbf{4.276} & 4 & GG & $<0.05$ \\
Average velocity & 2.631 & 4 & None & 0.0430 \\
Number of laps run & \textbf{8.771} & 4 & GG & $<0.001$ \\
Number of licks & \textbf{3.883} & 4 & GG & 0.0326 \\
Ratio of rewarded laps & 3.331 & 4 & None & 0.0157 \\
\bottomrule
\end{tabular}
\end{table}

\subsection{dHPC-to-NAc neurons are enriched for place cells with enhanced spatial tuning}

We used dual-color two-photon calcium imaging in Thy1-GCaMP6s mice to simultaneously record GCaMP6s activity while identifying NAc-projecting neurons via retrograde mCherry labeling (Fig.~2a--c). Across 6 mice and 19 imaging sessions we recorded 5,372 GCaMP-expressing neurons, of which 444 were identified as putative dHPC-to-NAc projection neurons and 4,928 as non-projecting (dHPC$^{-}$) neurons.

Place cells were identified by comparing spatial information content against 1,000 randomly shuffled distributions (95th percentile threshold). A significantly higher proportion of dHPC-to-NAc neurons were classified as place cells compared to dHPC$^{-}$ neurons (169/444, 38\% vs.\ 1,581/4,928, 32\%; $\chi^{2}(1, 5372)=6.364$, $p=0.012$; Table~\ref{tab:placecells}). Moreover, place cells in the projection population exhibited superior spatial tuning across all four metrics assessed by Welch's $t$-test (Table~\ref{tab:spatialtuning}).

\begin{table}[H]
\centering
\caption{Place cell proportions in dHPC$^{-}$ versus dHPC-to-NAc populations ($n=5{,}372$ neurons; 6 mice, 19 sessions). $\uparrow$ indicates metrics where higher values reflect better spatial coding.}
\label{tab:placecells}
\begin{tabular}{lcccc}
\toprule
Population & Place cells / Total & Proportion ($\uparrow$) & $\chi^{2}(1, 5372)$ & $p$ \\
\midrule
dHPC$^{-}$ & 1,581 / 4,928 & 32\% & \multirow{2}{*}{\textbf{6.364}} & \multirow{2}{*}{\textbf{0.012}} \\
dHPC-to-NAc & 169 / 444 & \textbf{38\%} & & \\
\bottomrule
\end{tabular}
\end{table}

\begin{table}[H]
\centering
\caption{Spatial tuning characteristics of place cells: dHPC$^{-}$ vs.\ dHPC-to-NAc (Welch's $t$-test; $n_{\text{dHPC}^{-}}=1{,}581$, $n_{\text{dHPC-to-NAc}}=169$ place cells). $\uparrow$ indicates metrics where higher values reflect better spatial coding.}
\label{tab:spatialtuning}
\begin{tabular}{lccc}
\toprule
Measure ($\uparrow$) & Welch's $t$ (df) & $p$ \\
\midrule
Spatial information rate & $t(186.55)=\mathbf{4.770}$ & $<0.001$ \\
Sparsity & $t(218.79)=\mathbf{3.657}$ & $<0.001$ \\
Reliability (in-field per lap) & $t(203.46)=\mathbf{2.479}$ & 0.014 \\
In-place field activity & $t(190.32)=\mathbf{2.884}$ & 0.0044 \\
\bottomrule
\end{tabular}
\end{table}

Place field width distributions also differed between populations (Kolmogorov--Smirnov test, $D=0.116$, $p=0.0299$), with dHPC-to-NAc place fields showing broader median width (58.9\,cm vs.\ 51.7\,cm; quartiles: 36.5/58.9/69.0\,cm vs.\ 35.6/51.7/69.0\,cm for dHPC$^{-}$).

\subsection{Place fields are anchored to local cues and enriched near the reward zone in dHPC-to-NAc neurons}

We examined whether place field boundaries were non-uniformly distributed and whether they differed between populations. Place field start and end positions were overrepresented at texture boundaries for both populations ($>$99.9th percentile compared to shuffled distributions). Critically, dHPC-to-NAc neurons showed significantly more place fields starting and ending near texture boundaries than dHPC$^{-}$ neurons ($\chi^{2}(1, 5134)=5.735$, $p=0.033$ for start positions; $\chi^{2}(1, 5217)=4.397$, $p=0.036$ for end positions; Table~\ref{tab:boundaries}).

\begin{table}[H]
\centering
\caption{Place field boundary enrichment near texture landmarks and reward zone: chi-squared comparisons of dHPC$^{-}$ vs.\ dHPC-to-NAc ($n=5{,}372$ neurons; 6 mice).}
\label{tab:boundaries}
\begin{tabular}{lccc}
\toprule
Comparison & $\chi^{2}$ (df) & $p$ \\
\midrule
PF start positions (dHPC-to-NAc vs.\ dHPC$^{-}$) & $\chi^{2}(1, 5134)=\mathbf{5.735}$ & 0.033 \\
PF end positions (dHPC-to-NAc vs.\ dHPC$^{-}$) & $\chi^{2}(1, 5217)=\mathbf{4.397}$ & 0.036 \\
PF center positions (texture boundaries) & $\chi^{2}(1, 5372)=1.646$ & 0.1995 \\
\bottomrule
\end{tabular}
\end{table}

To test whether reward zone representation depended on behavioral success and projection identity, we performed a two-way ANOVA with factors success (high vs.\ low) and projection (dHPC$^{-}$ vs.\ dHPC-to-NAc). This revealed a significant main effect of success ($F_{\text{success}}(1,1)=54.918$, $p<0.001$) but no main effect of projection ($F_{\text{projection}}(1,1)=0.958$, $p=0.338$) and a trending interaction ($F_{\text{interaction}}(1,1)=2.969$, $p=0.098$). Post-hoc Welch's $t$-tests with Bonferroni correction showed that dHPC-to-NAc neurons exhibited a robust reward zone overrepresentation on high-success sessions ($t(3.479)=8.671$, $p=0.0036$), whereas dHPC$^{-}$ neurons showed only a marginal effect ($t(3.561)=3.698$, $p=0.051$; Table~\ref{tab:rewardzone}).

\begin{table}[H]
\centering
\caption{Reward zone overrepresentation: two-way ANOVA and post-hoc Welch's $t$-tests with Bonferroni correction ($n=6$ mice, 19 sessions). $\uparrow$ indicates that higher overrepresentation is the expected direction for successful reward learning.}
\label{tab:rewardzone}
\begin{tabular}{lcccc}
\toprule
\multicolumn{5}{l}{\textit{Two-way ANOVA}} \\
Factor & $F$-statistic & df & $p$ \\
\midrule
Success & \textbf{54.918} & 1,1 & $<0.001$ \\
Projection & 0.958 & 1,1 & 0.338 \\
Interaction & 2.969 & 1,1 & 0.098 \\
\midrule
\multicolumn{5}{l}{\textit{Post-hoc Welch's $t$-tests (Bonferroni corrected)}} \\
Comparison ($\uparrow$) & $t$ (df) & $p$ \\
\midrule
dHPC$^{-}$ (high vs.\ low success) & $t(3.561)=3.698$ & 0.051 \\
dHPC-to-NAc (high vs.\ low success) & $t(3.479)=\mathbf{8.671}$ & \textbf{0.0036} \\
Low success (dHPC$^{-}$ vs.\ dHPC-to-NAc) & $t(13.629)=0.093$ & 1 \\
High success (dHPC$^{-}$ vs.\ dHPC-to-NAc) & $t(3.952)=2.075$ & 0.215 \\
\bottomrule
\end{tabular}
\end{table}

Behavioral success rate positively correlated with place field density near the reward zone for dHPC$^{-}$ neurons (Pearson $r(15)=0.622$, $p=0.0107$). Reward zone decoding using an SVM-based linear classifier (cross-validated on odd/even laps) confirmed that dHPC-to-NAc populations decoded reward zone presence significantly better than sample-size-matched dHPC$^{-}$ populations (Wilcoxon's $W(9)=5.0$, $n=10$ imaging sessions, $p=0.0195$).

\subsection{Speed-inhibited neurons are overrepresented in dHPC-to-NAc projections}

Velocity modulation was assessed by linear regression of calcium activity against 1\,cm/s velocity bins from 2 to 30\,cm/s, with Benjamini--Hochberg FDR correction at $p<0.05$~\citep{benjamini1995}. Speed-inhibited neurons were significantly overrepresented in the dHPC-to-NAc population (21\% vs.\ 15\%; $\chi^{2}(1, 5372)=13.66$, $p=0.00022$), whereas the proportion of speed-excited neurons did not differ significantly (16\% vs.\ 13\%; $\chi^{2}(1, 5372)=2.565$, $p=0.109$; Table~\ref{tab:velocity}).

\begin{table}[H]
\centering
\caption{Velocity modulation in dHPC$^{-}$ vs.\ dHPC-to-NAc populations (chi-squared test; $n=5{,}372$ neurons; 6 mice). $\downarrow$ for speed-inhibited indicates that lower velocity corresponds to higher firing. Bold indicates statistically significant differences.}
\label{tab:velocity}
\begin{tabular}{lccccc}
\toprule
Category & dHPC$^{-}$ (\%) & dHPC-to-NAc (\%) & $\chi^{2}(1, 5372)$ & $p$ \\
\midrule
Speed-excited & 13\% & 16\% & 2.565 & 0.109 \\
Speed-inhibited ($\downarrow$) & 15\% & \textbf{21\%} & \textbf{13.66} & \textbf{0.00022} \\
\bottomrule
\end{tabular}
\end{table}

Representative single-neuron examples confirmed robust velocity tuning. A speed-excited neuron showed a positive linear relationship (slope $= 7.43 \times 10^{-4}$, intercept $= -2.097 \times 10^{-3}$, $r=0.937$, $p=0.0058$), while a speed-inhibited neuron showed a negative relationship (slope $= -1.0437 \times 10^{-4}$, intercept $= 2.814 \times 10^{-3}$, $r=-0.933$, $p=0.0065$).

\subsection{Appetitive lick-excited neurons are enriched in dHPC-to-NAc projections}

We examined lick-related modulation using a two-way mixed ANOVA on population-level activity with factors lick timing (pre vs.\ post) and projection (dHPC$^{-}$ vs.\ dHPC-to-NAc). This revealed a significant main effect of projection ($F_{\text{projection}}(1, 5370)=43.779$, $p<0.001$) and a significant interaction ($F_{\text{interaction}}(1, 5370)=7.073$, $p=0.0079$), but no main effect of lick timing ($F_{\text{lick timing}}(1, 5370)=2.843$, $p=0.0918$). Post-hoc $t$-tests with Bonferroni correction confirmed that dHPC-to-NAc neurons showed significant lick-related modulation ($t(443)=2.470$, $p=0.0277$) while dHPC$^{-}$ neurons did not ($t(4927)=0.871$, $p=0.768$).

At the single-cell level, appetitive lick-excited neurons were significantly overrepresented among dHPC-to-NAc neurons (7.4\% vs.\ 3.8\%; $\chi^{2}(1, 5372)=13.018$, $p=0.00031$), while the proportion of lick-inhibited neurons was comparable (17.1\% vs.\ 19.7\%; $\chi^{2}(1, 5372)=1.626$, $p=0.202$; Table~\ref{tab:licking}). Total lick-modulated neurons comprised 1,268 of 5,372 (24\%).

\begin{table}[H]
\centering
\caption{Appetitive lick modulation: population ANOVA and single-cell chi-squared tests ($n=5{,}372$ neurons; 6 mice, 19 sessions). $\uparrow$ indicates metrics where higher values reflect stronger appetitive coding.}
\label{tab:licking}
\begin{tabular}{lccccc}
\toprule
\multicolumn{6}{l}{\textit{Two-way mixed ANOVA}} \\
Factor & $F$-statistic & df & $p$ \\
\midrule
Lick timing & 2.843 & 1, 5370 & 0.0918 \\
Projection & \textbf{43.779} & 1, 5370 & $<0.001$ \\
Interaction & \textbf{7.073} & 1, 5370 & 0.0079 \\
\midrule
\multicolumn{6}{l}{\textit{Single-cell proportions (chi-squared test)}} \\
Cell type ($\uparrow$) & dHPC$^{-}$ (\%) & dHPC-to-NAc (\%) & $\chi^{2}(1, 5372)$ & $p$ \\
\midrule
Lick-excited & 3.8\% & \textbf{7.4\%} & \textbf{13.018} & \textbf{0.00031} \\
Lick-inhibited & 19.7\% & 17.1\% & 1.626 & 0.202 \\
\bottomrule
\end{tabular}
\end{table}

Supplementary analyses confirmed these findings. Consummatory licking depression was analyzed by two-way mixed ANOVA, revealing a significant effect of reward timing ($F(2, 10740)=19.380$, Greenhouse--Geisser corrected, $p<0.001$) and a significant interaction ($F(2, 10740)=5.559$, $p=0.0039$). Lick modulation indices were significantly larger for dHPC-to-NAc neurons (pre-post difference: Welch's $t(496.11)=2.379$, $p=0.0177$; lick index: Welch's $t(454.19)=2.729$, $p=0.0066$).

\subsection{Optogenetic activation of dHPC terminals in NAc evokes appetitive motor output}

To test the causal significance of the dHPC-to-NAc pathway, we injected AAV2-CaMKII-hChR2(H134R)-EYFP-WPRE ($4 \times 10^{12}$ TU) bilaterally into dHPC and implanted bilateral fiber-optic cannulas over NAc ($n=4$ ChR2 mice, $n=3$ EYFP control mice). Blue light stimulation (473\,nm, 5\,mW, 20\,Hz, 5\,ms pulses) of dHPC terminals in NAc induced two key behavioral effects (Table~\ref{tab:opto}): (1) mouth movement, confirmed by paired $t$-test (ChR2: $t(3)=7.485$, $p=0.00494$; EYFP: $t(2)=1.353$, $p=0.309$), and (2) locomotor deceleration (ChR2: $t(3)=-3.551$, $p=0.0381$; EYFP: $t(2)=-1.263$, $p=0.334$). Neither effect was observed in EYFP controls.

\begin{table}[H]
\centering
\caption{Optogenetic stimulation of dHPC terminals in NAc: paired $t$-tests ($n=4$ ChR2 mice, $n=3$ EYFP controls). $\uparrow$ mouth motion and $\downarrow$ velocity indicate appetitive-like responses.}
\label{tab:opto}
\begin{tabular}{lcccc}
\toprule
Behavioral measure & Group & $t$ (df) & $p$ \\
\midrule
Mouth motion ($\uparrow$) & ChR2 & $t(3)=\mathbf{7.485}$ & \textbf{0.00494} \\
 & EYFP & $t(2)=1.353$ & 0.309 \\
Velocity change ($\downarrow$) & ChR2 & $t(3)=\mathbf{-3.551}$ & \textbf{0.0381} \\
 & EYFP & $t(2)=-1.263$ & 0.334 \\
\bottomrule
\end{tabular}
\end{table}

\subsection{Conjunctive coding is enriched in dHPC-to-NAc neurons}

We asked whether individual dHPC-to-NAc neurons encode conjunctions of spatial, velocity, and licking signals beyond what is predicted from single-feature tuning alone. Pairwise chi-squared analyses revealed significant enrichment of space $\times$ velocity conjunctive coding in the dHPC-to-NAc population. Among place cells, speed-inhibited neurons were overrepresented in dHPC-to-NAc relative to dHPC$^{-}$ ($\chi^{2}(1, 1750)=8.977$, $p=0.0027$). Speed-excited non-place cells were also enriched in dHPC-to-NAc ($\chi^{2}(1, 3622)=6.255$, $p<0.05$).

For space $\times$ licking conjunctions, lick-excited place cells were overrepresented in dHPC-to-NAc neurons ($\chi^{2}(1, 1750)=9.442$, $p<0.05$), with 15\% of dHPC-to-NAc place cells being lick-excited compared to 8\% of dHPC$^{-}$ place cells ($p<0.001$). Among non-place cells, lick-excited proportions did not differ ($\chi^{2}(1, 3622)=0.488$, $p=0.485$).

For velocity $\times$ licking conjunctions, speed-inhibited lick-excited neurons were enriched in dHPC-to-NAc relative to dHPC$^{-}$ ($\chi^{2}(1, 830)=6.564$, $p<0.05$), with 20\% of dHPC-to-NAc speed-inhibited neurons being lick-excited compared to 11\% in dHPC$^{-}$ ($p<0.001$).

\subsection{GLM confirms multimodal conjunctive coding in dHPC-to-NAc neurons}

To formally quantify conjunctive coding, we fitted a generalized linear Poisson model (GLM) using position, velocity, and appetitive licking as predictors, achieving approximately 40\% variance explained on held-out test data (Table~\ref{tab:glm}). Chi-squared tests on the proportions of neurons significantly modulated by each feature confirmed that dHPC-to-NAc neurons were enriched for all three single features: position ($\chi^{2}(1, 5372)=93.634$, $p<0.001$), velocity ($\chi^{2}(1, 5372)=141.86$, $p<0.001$), and licking ($\chi^{2}(1, 5372)=10.050$, $p=0.0015$).

\begin{table}[H]
\centering
\caption{GLM-based conjunctive coding: chi-squared tests comparing dHPC$^{-}$ vs.\ dHPC-to-NAc proportions ($n=5{,}372$ neurons; 6 mice). $\uparrow$ indicates that higher proportions reflect richer multimodal coding in the projection population. Bold indicates statistical significance at $p<0.05$.}
\label{tab:glm}
\begin{tabular}{lccc}
\toprule
Feature / Combination ($\uparrow$) & $\chi^{2}(1, 5372)$ & $p$ \\
\midrule
\multicolumn{3}{l}{\textit{Single features}} \\
Position & \textbf{93.634} & $<0.001$ \\
Velocity & \textbf{141.86} & $<0.001$ \\
Licking & \textbf{10.050} & \textbf{0.0015} \\
\midrule
\multicolumn{3}{l}{\textit{Conjunctive combinations}} \\
Position \& Velocity & \textbf{163.97} & $<0.001$ \\
Position \& Licking & \textbf{26.029} & $<0.001$ \\
Velocity \& Licking & \textbf{27.145} & $<0.001$ \\
Position \& Velocity \& Licking & \textbf{34.993} & $<0.001$ \\
\midrule
\multicolumn{3}{l}{\textit{N-feature coding categories}} \\
Non-coding (overrepresented in dHPC$^{-}$) & \textbf{35.382} & $<0.001$ \\
Single-coding (comparable) & 0.0057 & $>0.05$ \\
Dual-coding ($\uparrow$ dHPC-to-NAc) & \textbf{61.336} & $<0.001$ \\
Triple-coding ($\uparrow$ dHPC-to-NAc) & \textbf{30.447} & $<0.001$ \\
\bottomrule
\end{tabular}
\end{table}

Overall, 44\% of dHPC-to-NAc neurons exhibited conjunctive coding (responding to two or more features) compared to 19\% of dHPC$^{-}$ neurons ($p<0.001$, chi-squared). Linear decoding of reward zone presence using an SVM classifier showed that conjunctive coding neurons provided significantly more accurate reward zone predictions than non-conjunctive coding neurons (Wilcoxon's $W(18)=14.0$, $p<0.001$).

Supplementary GLM analyses revealed higher overall model fit for dHPC-to-NAc neurons (Welch's $t(615.15)=2.147$, $p=0.032$). Feature importance analysis showed significantly greater contributions of position (Welch's $t(2508.09)=2.162$, $p=0.031$) and licking (Welch's $t(1879.75)=9.328$, $p<0.001$) but not velocity (Welch's $t(2609.15)=0.881$, $p=0.378$) in dHPC-to-NAc neurons. Feature ablation ($R^{2}$ drop) confirmed stronger dependence on position (Welch's $t(523.23)=5.984$, $p<0.001$) and velocity (Welch's $t(497.69)=7.704$, $p<0.001$), but not licking (Welch's $t(552.63)=0.3626$, $p=0.717$; Table~\ref{tab:glmsup}).

\begin{table}[H]
\centering
\caption{Supplementary GLM analysis: feature importance and $R^{2}$ drop (Welch's $t$-test; dHPC$^{-}$ vs.\ dHPC-to-NAc; $n=5{,}372$ neurons). $\uparrow$ indicates that higher values reflect greater feature contribution. Bold indicates $p<0.05$.}
\label{tab:glmsup}
\begin{tabular}{lccc}
\toprule
Measure ($\uparrow$) & Welch's $t$ (df) & $p$ \\
\midrule
\multicolumn{3}{l}{\textit{Overall model fit}} \\
GLM $R^{2}$ & $t(615.15)=\mathbf{2.147}$ & \textbf{0.032} \\
\midrule
\multicolumn{3}{l}{\textit{Feature importance}} \\
Position & $t(2508.09)=\mathbf{2.162}$ & \textbf{0.031} \\
Licking & $t(1879.75)=\mathbf{9.328}$ & $<0.001$ \\
Velocity & $t(2609.15)=0.881$ & 0.378 \\
\midrule
\multicolumn{3}{l}{\textit{$R^{2}$ drop (feature ablation)}} \\
Position & $t(523.23)=\mathbf{5.984}$ & $<0.001$ \\
Velocity & $t(497.69)=\mathbf{7.704}$ & $<0.001$ \\
Licking & $t(552.63)=0.3626$ & 0.717 \\
\bottomrule
\end{tabular}
\end{table}

% ============================================================
% DISCUSSION
% ============================================================
\section{Discussion}

Our results demonstrate that the dorsal hippocampus routes an enriched, multimodal code to the nucleus accumbens that integrates spatial, locomotor, and appetitive information to support goal-directed behavior. By combining dual-color two-photon imaging with retrograde labeling to record 5,372 hippocampal neurons including 444 identified NAc-projecting cells, we provide the most comprehensive characterization to date of projection-specific coding in the hippocampal--accumbal pathway during appetitive navigation.

\paragraph{Enhanced spatial coding in hippocampal--accumbal projections.}
We found that dHPC-to-NAc neurons are enriched for place cells (38\% vs.\ 32\%; $p=0.012$) and display superior spatial tuning on all four metrics assessed: spatial information rate ($p<0.001$), sparsity ($p<0.001$), reliability ($p=0.014$), and in-field activity ($p=0.0044$). These findings extend prior work on projection-specific hippocampal coding~\citep{ciocchi2015, danielson2016} by demonstrating that the NAc-projecting subpopulation carries a sharper, more informative spatial signal than the general hippocampal population. The slightly broader place fields in dHPC-to-NAc neurons (median 58.9\,cm vs.\ 51.7\,cm) may reflect a coding strategy that favors spatial coverage near behaviorally relevant landmarks while maintaining high information content per spike~\citep{skaggs1996, sosa2018}.

\paragraph{Reward zone enrichment and behavioral relevance.}
Place fields in dHPC-to-NAc neurons were preferentially anchored to texture boundaries and showed robust overrepresentation near the reward zone during high-success sessions ($t(3.479)=8.671$, $p=0.0036$), an effect that was only marginal in non-projecting neurons ($p=0.051$). This pattern suggests that the hippocampal--accumbal pathway selectively amplifies reward-relevant spatial representations when animals successfully engage in goal-directed behavior. The correlation between behavioral success and reward zone place field density ($r=0.622$, $p=0.0107$) further supports the functional significance of this enrichment~\citep{dupret2010, gauthier2018}. Critically, SVM-based decoding confirmed that dHPC-to-NAc populations carry superior reward zone information compared to sample-matched controls ($W(9)=5.0$, $p=0.0195$).

\paragraph{Velocity and appetitive signals converge in the projection pathway.}
Beyond spatial coding, dHPC-to-NAc neurons exhibited enriched velocity and lick-related modulation. The overrepresentation of speed-inhibited neurons (21\% vs.\ 15\%; $p=0.00022$) suggests that the projection pathway preferentially signals deceleration -- a behavioral correlate of reward approach and consumption~\citep{markus1995}. Appetitive lick-excited neurons were nearly twice as prevalent in the projection population (7.4\% vs.\ 3.8\%; $p=0.00031$), indicating that the hippocampus actively transmits consummatory motor-related signals to the NAc. These findings align with the hypothesis that the NAc integrates spatial and motivational information from hippocampal inputs to coordinate approach behavior~\citep{mogenson1980, pennartz2011, floresco2015}.

\paragraph{Causal evidence from optogenetic activation.}
Optogenetic stimulation of dHPC terminals in NAc evoked two complementary appetitive-like responses: orofacial movement ($t(3)=7.485$, $p=0.00494$) and locomotor deceleration ($t(3)=-3.551$, $p=0.0381$), neither of which was observed in EYFP controls. These results provide direct causal evidence that activation of the hippocampal--accumbal pathway is sufficient to trigger appetitive motor output, bridging the gap between correlational coding analyses and behavioral function~\citep{boyden2005, britt2012, leshan2014}.

\paragraph{Conjunctive multimodal coding supports reward zone representation.}
Perhaps the most striking finding is the dramatic enrichment of conjunctive coding in dHPC-to-NAc neurons. GLM analysis revealed that 44\% of projection neurons encoded conjunctions of position, velocity, and licking signals, compared to only 19\% of non-projecting neurons ($p<0.001$). This multimodal conjunction was not simply a byproduct of single-feature enrichment, as formal comparisons of dual- and triple-coding categories all exceeded expectation ($\chi^{2}$ values of 61.336 and 30.447, respectively; both $p<0.001$). Critically, conjunctive coding neurons supported more accurate reward zone decoding than non-conjunctive neurons ($W(18)=14.0$, $p<0.001$), suggesting that multimodal integration at the single-neuron level provides a computational advantage for representing goal-relevant spatial information~\citep{eichenbaum2014, trouche2019}.

\paragraph{Implications and limitations.}
Our findings redefine the hippocampal--accumbal pathway as a multimodal information channel that transmits not just spatial but also locomotor and appetitive signals in a conjunctive format optimized for goal-directed behavior. This view extends the classic framework of the NAc as a limbic--motor interface~\citep{mogenson1980, groenewegen2007} by specifying the computational content of the hippocampal input stream. Future work should determine whether similar projection-specific coding principles apply to other hippocampal output pathways (e.g., to prefrontal cortex or lateral septum) and whether the enriched conjunctive code is shaped by learning~\citep{dupret2010}.

Several limitations should be noted. First, calcium imaging provides an indirect measure of spiking activity with limited temporal resolution compared to electrophysiology~\citep{chen2013}. Second, our retrograde labeling strategy identifies neurons with axons in NAc but cannot distinguish between shell and core subregions. Third, optogenetic experiments used relatively small sample sizes ($n=4$ ChR2, $n=3$ EYFP) and terminal stimulation may produce supraphysiological activation patterns. Fourth, all recordings were performed in head-fixed mice on a linear belt, which constrains the spatial complexity of the navigation environment~\citep{dombeck2010}.

% ============================================================
% METHODS
% ============================================================
\section{Methods}

\subsection{Animals}
All procedures were approved by the local animal care committee and conducted in accordance with institutional guidelines. We used Thy1-GCaMP6s transgenic mice~\citep{dana2014} for calcium imaging experiments ($n=6$ mice, 19 imaging sessions) and C57BL/6 mice for optogenetic experiments ($n=4$ ChR2, $n=3$ EYFP). A total of $n=18$ mice (two cohorts of 9) were used for behavioral training.

\subsection{Viral injections and surgical procedures}
Mice were anesthetized with ketamine (0.13\,mg/g) and xylazine (0.01\,mg/g, i.p.) and received post-surgical analgesia with buprenorphine (0.05\,mg/kg, administered thrice daily for 3 days). For retrograde labeling, AAVrg-pgk-Cre (Addgene \#24593; titer: $1 \times 10^{13}$\,vg/ml) was injected into the ipsilateral right NAc ($-1.3$\,mm AP, $-1.0$\,mm lateral, 5.0 and 4.3\,mm ventral relative to Bregma; $2 \times 500$\,nl at 100\,nl/min). AAV5-hSyn-DIO-mCherry (Addgene \#50459; titer: $1.1 \times 10^{13}$\,vg/ml; 200\,nl) was injected into ipsilateral dHPC (3.38\,mm AP, $-2.5$\,mm lateral, 1.8\,mm ventral; 10-degree angle). A stainless-steel cranial window cannula (1.7\,mm long, 3\,mm outer diameter) with a 3\,mm coverslip was implanted over dHPC through a 3\,mm craniotomy.

For optogenetic experiments, AAV2-CaMKII-hChR2(H134R)-EYFP-WPRE (UNC \#AV4381E; titer: $4 \times 10^{12}$\,TU) or rAAV2-CaMKII-EYFP (UNC \#AV6650; titer: $4.3 \times 10^{12}$\,TU) was injected bilaterally into dHPC (200\,nl each; 3.38\,mm AP, $\pm$2.5\,mm lateral, 1.8\,mm ventral; 10-degree angle). Bilateral fiber-optic cannulas were implanted over NAc ($+1.3$\,mm AP, $\pm 1$\,mm lateral, $-4.9$\,mm ventral from brain surface)~\citep{boyden2005}.

\subsection{Two-photon calcium imaging}
Dual-color two-photon imaging was performed using a Ti:sapphire laser (Chameleon Ultra II) at 920\,nm for GCaMP6s excitation and a fixed-wavelength laser (Spectra Physics) at 1045\,nm for mCherry excitation, through a 16$\times$ objective (N16XLWD-PF, Nikon). Emission was collected through 525/40\,nm (GCaMP) and 590/40\,nm (mCherry) bandpass filters. Images were acquired at 15.2\,Hz ($1024 \times 1024$ pixels) or 30.5\,Hz ($512 \times 512$ pixels), with fields of view spanning 350--850\,$\mu$m. Signal collection rate was 10\,kHz.

\subsection{Behavioral task}
Mice were food-restricted to 80\% of daily food pellets (10--20\% weight loss) and trained for 5 days (15\,min/day) on a 360\,cm textile belt (7\,cm width) with 6 textured zones, a 30\,cm hidden reward zone, and a 30\,cm anticipation zone. Behavior was monitored by camera at 25 or 75\,Hz ($782 \times 582$ pixels). Successful laps were defined as those with $\geq$50\% relative licking in the reward plus anticipation zones with reward received; high-success sessions had $\geq$50\% successful laps.

\subsection{Calcium imaging analysis}
Motion correction was performed using NormCorre~\citep{pnevmatikakis2016} (max\_shifts $= 40$, num\_frames\_split $= 2000$, overlaps $= 46$, splits\_els $= 4$, strides $= 255$). Cell detection and signal extraction used CaImAn v1.6.2~\citep{giovannucci2019}. Spatial activity maps were computed in 45 bins of 8\,cm for periods with velocity $>$2\,cm/s.

\subsection{Place cell identification}
Spatial information content was compared against 1,000 randomly shuffled distributions (circular shifts of spike times). Neurons exceeding the 95th percentile of shuffled spatial information were classified as place cells~\citep{skaggs1996}.

\subsection{Velocity modulation analysis}
Calcium activity was binned in 1\,cm/s velocity bins from 2 to 30\,cm/s and assessed by linear regression. Significance was determined at $p<0.05$ after Benjamini--Hochberg FDR correction~\citep{benjamini1995}. Neurons with significant positive slopes were classified as speed-excited; those with significant negative slopes as speed-inhibited.

\subsection{Appetitive lick modulation}
Appetitive lick bouts were defined by $<$2\,s between events, with onset preceded by $\geq$3\,s absence of licks. Pre/post calcium activity was compared using Wilcoxon's non-parametric paired test. Population-level analysis used two-way mixed ANOVA (factors: lick timing, projection identity) followed by post-hoc $t$-tests with Bonferroni correction.

\subsection{Generalized linear model}
A Poisson GLM with position, velocity, and appetitive licking as predictors was trained (80\% train, 10\% validation, 10\% test; Gaussian temporal smoothing, $\sigma = 2$ frames, 0.5\,s window; downsampled to 3\,Hz). Significance was determined by comparison against $3 \times 100$ shuffled models per feature; a neuron was classified as significantly modulated if its true $R^{2}$ exceeded 96 of 100 shuffled models ($>$95\%).

\subsection{Reward zone decoding}
An SVM-based linear classifier was cross-validated on odd/even laps to decode the presence of the reward zone from population activity. Population comparisons used sample-size-matched random subsets of the larger population.

\subsection{Optogenetic stimulation}
Blue light (473\,nm) was delivered at 5\,mW, 20\,Hz (5\,ms pulses) for up to 10\,s through bilateral fiber-optic cannulas. Behavioral effects (mouth motion, velocity change) were assessed by paired $t$-tests comparing stimulation vs.\ baseline periods.

\subsection{Statistical analysis}
All data are presented as mean $\pm$ SEM. Significance was set at $p<0.05$. The following tests were used: repeated-measures ANOVA (with Greenhouse--Geisser correction where sphericity was violated), chi-squared test, Welch's $t$-test, two-way ANOVA, two-way mixed ANOVA, post-hoc $t$-tests with Bonferroni correction, Wilcoxon's signed-rank test, paired $t$-test, permutation test (1,000$\times$ bootstrap), Kolmogorov--Smirnov test, Pearson correlation, linear regression, and generalized linear Poisson model with shuffle-based significance testing.

% ============================================================
% CONCLUSION
% ============================================================
\section{Conclusion}

We have demonstrated that dorsal hippocampal neurons projecting to the nucleus accumbens carry an enriched, multimodal code that integrates spatial position, locomotor velocity, and appetitive licking signals. This projection population is enriched for place cells with sharper spatial tuning, preferentially encodes reward zone location during successful behavior, and shows dramatically elevated conjunctive coding (44\% vs.\ 19\%). Optogenetic activation of this pathway evokes appetitive motor output including orofacial movement and deceleration. These findings establish the hippocampal--accumbal projection as a dedicated channel for transmitting behaviorally relevant, conjunctive spatial--motivational information that guides goal-directed appetitive navigation.

\section*{Data availability}
The data that support the findings of this study are available from the corresponding author upon reasonable request.

\section*{Acknowledgements}
We thank all members of the laboratory for discussions and technical assistance.

\bibliographystyle{unsrtnat}
\bibliography{references}

\end{document}
